
\documentclass{article}
\usepackage{colortbl}
\usepackage{makecell}
\usepackage{multirow}
\usepackage{supertabular}

\begin{document}

\newcounter{utterance}

\centering \large Interaction Transcript for game `cladder', experiment `full\_v1.5\_default', episode 4490 with Qwen3.6{-}27B.
\vspace{24pt}

{ \footnotesize  \setcounter{utterance}{1}
\setlength{\tabcolsep}{0pt}
\begin{supertabular}{c@{$\;$}|p{.15\linewidth}@{}p{.15\linewidth}p{.15\linewidth}p{.15\linewidth}p{.15\linewidth}p{.15\linewidth}}
    \# & \multicolumn{2}{c}{Player} && \multicolumn{2}{c}{Game Master} \\
    \hline

    \theutterance \stepcounter{utterance}  
    & & & \multicolumn{4}{p{0.6\linewidth}}{
        \cellcolor[rgb]{0.9,0.9,0.9}{
            \makecell[{{p{\linewidth}}}]{
                \texttt{\tiny{[P1$\langle$GM]}}
                \texttt{You are an expert in causal inference. The following question is not a typical commonsense query, but rather a meticulously designed question created by a professor specializing in causal inference, intended to assess the students' mastery of the course content.} \\
\texttt{Imagine a self{-}contained, hypothetical world with only the following conditions, and without any unmentioned factors or causal relationships: Listening to jazz has a direct effect on tar deposit. Tar deposit has a direct effect on lung cancer. For people who do not listen to jazz, the probability of lung cancer is 61\%. For people who listen to jazz, the probability of lung cancer is 79\%. For people who listen to jazz, would it be less likely to see lung cancer if the person did not listen to jazz?} \\
\texttt{Start your answer with 'Yes' or 'No', followed by additional reasoning or evidence to support your explanation.} \\
            }
        }
    }
     \\ \\

    \theutterance \stepcounter{utterance}  
    & \multicolumn{4}{p{0.6\linewidth}}{
        \cellcolor[rgb]{0.561,0.737,0.561}{
            \makecell[{{p{\linewidth}}}]{
                \texttt{\tiny{[P1$\rangle$GM]}}
                \texttt{No} \\
\\ 
\texttt{**Reasoning:**} \\
\\ 
\texttt{To answer this question, we must distinguish between the **observed association** (correlation) and the **counterfactual causal effect** for a specific individual.} \\
\\ 
\texttt{1.  **Identify the Causal Structure:**} \\
\texttt{    The problem states:} \\
\texttt{    *   Jazz \$\textbackslash{}rightarrow\$ Tar Deposit \$\textbackslash{}rightarrow\$ Lung Cancer.} \\
\texttt{    *   There are no other factors or relationships.} \\
\texttt{    *   This implies that Jazz affects Lung Cancer *only* through Tar Deposit.} \\
\\ 
\texttt{2.  **Analyze the Observed Probabilities:**} \\
\texttt{    *   \$P(\textbackslash{}text\{Lung Cancer\} | \textbackslash{}text\{No Jazz\}) = 0.61\$} \\
\texttt{    *   \$P(\textbackslash{}text\{Lung Cancer\} | \textbackslash{}text\{Jazz\}) = 0.79\$} \\
\texttt{    *   The observed risk difference is \$0.79 {-} 0.61 = 0.18\$. This means that, on average, people who listen to jazz have an 18\% higher probability of lung cancer than those who do not.} \\
\\ 
\texttt{3.  **Analyze the Counterfactual Question:**} \\
\texttt{    The question asks: *"For people who listen to jazz, would it be less likely to see lung cancer if the person did not listen to jazz?"*} \\
\\ 
\texttt{    This is asking about the **Individual Causal Effect** (or the counterfactual outcome) for a specific person who *} \\
            }
        }
    }
    & & \\ \\

    \theutterance \stepcounter{utterance}  
    & & & \multicolumn{2}{p{0.3\linewidth}}{
        \cellcolor[rgb]{0.9,0.9,0.9}{
            \makecell[{{p{\linewidth}}}]{
                \texttt{\tiny{[GM$|$GM]}}
                \texttt{No} \\
            }
        }
    }
    & & \\ \\

    \theutterance \stepcounter{utterance}  
    & & & \multicolumn{2}{p{0.3\linewidth}}{
        \cellcolor[rgb]{0.9,0.9,0.9}{
            \makecell[{{p{\linewidth}}}]{
                \texttt{\tiny{[GM$|$GM]}}
                \texttt{yes} \\
            }
        }
    }
    & & \\ \\

    \theutterance \stepcounter{utterance}  
    & & & \multicolumn{2}{p{0.3\linewidth}}{
        \cellcolor[rgb]{0.9,0.9,0.9}{
            \makecell[{{p{\linewidth}}}]{
                \texttt{\tiny{[GM$|$GM]}}
                \texttt{game\_result = LOSE} \\
            }
        }
    }
    & & \\ \\

\end{supertabular}
}

\end{document}
